\documentclass[11pt]{article}

\usepackage[margin=1in]{geometry}
\usepackage{xcolor}
\usepackage{enumitem}
\usepackage{hyperref}
\usepackage{amsmath,amssymb}
\usepackage{lmodern}

\hypersetup{
  colorlinks=true,
  linkcolor=blue,
  citecolor=blue,
  urlcolor=blue
}

% ---------- Macros ----------
\newcommand{\manuscripttitle}{[Paper Title]}
\newcommand{\journalname}{[Journal / Conference Name]}
\newcommand{\manuscriptid}{[Manuscript ID]}
\newcommand{\corrauthor}{[Corresponding Author Name]}

\newcommand{\comment}[1]{%
  \vspace{0.6em}
  \noindent\textbf{Reviewer comment.} \emph{#1}
  \vspace{0.3em}
}

\newcommand{\response}[1]{%
  \noindent\textbf{Response.} #1
  \vspace{0.6em}
}

\newcommand{\revision}[1]{%
  \noindent\textbf{Revision made.} #1
  \vspace{0.6em}
}

\newcommand{\location}[1]{%
  \noindent\textbf{Location in revised manuscript.} #1
  \vspace{0.8em}
}

\begin{document}

% \begin{center}
%   {\Large \textbf{Response to Reviewers}}\\[0.5em]
%   \textbf{Manuscript Title:} \manuscripttitle\\
%   \textbf{Manuscript ID:} \manuscriptid\\
%   \textbf{Submitted to:} \journalname
% \end{center}

\vspace{1em}



\vspace{0.5em}


\newpage

% ============================================================
\section*{Response to Reviewer 1}

\comment{
1. Stationarity Relationship Between the Transformed Primal–Dual Problem and the Original Problem.
}

\response{
We thank the reviewer for raising this important concern.

Here, let us first define the problem and primal-dual function
\[
(P)
x,B
min
	​

J(x,B):=F(x,B)+r(B),
L(x,B,λ)=F(x,B)+⟨λ,B⟩−r
∗
(λ)
\]
Under the condition $r(x) = r**$, suppose $B$ is the primal-dual stationarity point of Problem (P), then 

 
we clarify the relationship between reduced primal--dual stationarity and stationarity of the original penalized problem and give the 

(1) 
Let $r$ be proper and lower semicontinuous, and suppose that $(x^*,B^*,\lambda^*)$ satisfies $\nabla_xF(x^*,B^*)=0$, $\nabla_BF(x^*,B^*)+\lambda^*=0$, and $B^*\in\partial r^*(\lambda^*)$. By subgradient inversion, $\lambda^*\in\partial r^{**}(B^*)$. If the contact condition $r(B^*)=r^{**}(B^*)<+\infty$ holds, the supporting-subgradient inequality for $r^{**}$, together with $r\ge r^{**}$, implies $\lambda^*\in\partial r(B^*)$. Consequently, $0\in\nabla_BF(x^*,B^*)+\partial r(B^*)$, and $(x^*,B^*)$ is a first-order stationary point of $F+r$. For the W-type regularizer $r(B)=\kappa\sum_{i,j}\bigl||B_{ij}|-1\bigr|$, the contact condition is equivalent to $|B_{ij}^*|\ge1$ for all $i,j$. Thus, reduced stationarity transfers to the original penalized problem only under this additional contact condition.
}

% \revision{
% We revised the discussion of [topic] and added [new explanation / theorem statement / proof paragraph / experimental result].
% }

% \location{
% Section~[X], page~[Y], lines~[A--B].
% }

% \comment{
% The reviewer asked whether [second comment].
% }

% \response{
% This is a helpful question. The answer is [direct answer]. The reason is [technical explanation]. In particular, [connect to theorem, assumption, experiment, or algorithmic detail].
% }

% \revision{
% We added a clarification to prevent this ambiguity.
% }

% \location{
% Section~[X], page~[Y], lines~[A--B].
% }

% % ============================================================
% \section*{Response to Reviewer 2}

% \comment{
% The reviewer commented that [quote or paraphrase].
% }

% \response{
% We appreciate the reviewer pointing this out. We agree that the original version did not explain [issue] sufficiently. We have now expanded the discussion and clarified [specific correction].

% The revised explanation emphasizes [main mechanism]. This is supported by [proof / proposition / experiment / table / figure].
% }

% \revision{
% We added [specific revision] and modified [specific sentence, paragraph, theorem, table, or figure].
% }

% \location{
% Section~[X], page~[Y], lines~[A--B].
% }

% \comment{
% The reviewer requested additional comparison with [method / paper / baseline].
% }

% \response{
% We thank the reviewer for this suggestion. We have added a comparison with [method / reference]. The key distinction is that our method [precise novelty claim], while [baseline / prior work] assumes or requires [precise limitation].

% This comparison has been added to make the contribution clearer.
% }

% \revision{
% We expanded the related work and added [citation / comparison paragraph / table entry].
% }

% \location{
% Section~[X], page~[Y], lines~[A--B].
% }

% % ============================================================
% \section*{Response to Reviewer 3}

% \comment{
% The reviewer noted that [quote or paraphrase].
% }

% \response{
% We thank the reviewer for identifying this issue. We have corrected [error / notation / unclear statement]. The correction does not affect the main results because [brief technical reason].
% }

% \revision{
% We corrected the notation and updated the surrounding explanation.
% }

% \location{
% Section~[X], page~[Y], lines~[A--B].
% }

% \comment{
% The reviewer suggested improving the presentation of [topic].
% }

% \response{
% We agree with this suggestion. We have reorganized the relevant part to make the logic clearer. The revised version now first introduces [concept], then states [result], and finally explains [implication].
% }

% \revision{
% We rewrote the paragraph on [topic] and added [figure / table / explanatory text] where appropriate.
% }

% \location{
% Section~[X], page~[Y], lines~[A--B].
% }

% % ============================================================
% \section*{Additional Changes}

% In addition to the reviewer-specific revisions above, we made the following minor changes:
% \begin{enumerate}[leftmargin=2em]
%   \item Corrected typographical errors in Sections~[X] and~[Y].
%   \item Unified notation for [symbol / quantity / operator].
%   \item Improved the caption of Figure~[X].
%   \item Added missing references to [citations].
% \end{enumerate}

% \vspace{1em}

% \noindent We again thank the editor and reviewers for their time and constructive feedback. We hope that the revised manuscript adequately addresses the comments and is now suitable for publication in \journalname.

% \vspace{1.5em}

% \noindent Sincerely,\\
% \corrauthor\\
% on behalf of all authors

\end{document} 
